Patients are often asked to prove a simple fact about their health record,
such as having received at least a required number of vaccine doses, and
end up disclosing the whole record. Zero-knowledge proofs can answer such a
question without revealing the data, but a proof is only meaningful if it is
tied to a stored record and to a verification session. We start from a
prototype that stored an encrypted record on IPFS and its content identifier
on Solana, in which a server computed the threshold proof from the private
value and the proof referred to no record. We redesign the proof path: a
salted Poseidon commitment to the private value is stored in the record's
on-chain account and is also an output of the circuit; a verifier-chosen
challenge is a public input; input ranges are enforced inside the circuit;
and proofs are generated in the owner's browser. We implement the design
with Circom, snarkjs and the Anchor framework and measure it on a laptop.
The bound circuit has \NewConstraints{} constraints against
\OldConstraints{} for the baseline. A Groth16 proof takes a median of
\BrowserProveMs~ms to generate in a browser and \NewVerifyMs~ms to verify.
Each record occupies a \AccountBytes-byte account with a refundable deposit
of \RentSol{} SOL, and its creation consumes at most \CuMax{} compute units.
We analyse the security of the design and state its limitations, chief among
them that the committed value is self-asserted.
